\subsection{Springer 纤维的例子}

\subsubsection{} 当 $e=0$ 时，$\sB_{e}=\sB$。

\subsubsection{正则幂零元}\label{sss:regular Spr} 唯一的稠密开 $G$ 轨道由{\em 正则幂零元}组成，即满足 $\dim C_{G}(e)=r$ 的那些 $e$。当 $e$ 是正则幂零元时，$\sB_{e}$ 是单点：恰有一个 Borel 子群的 Lie 代数包含 $e$。这个 Borel 子群是什么？

由 Jacobson--Morosov 定理，可以把 $e$ 扩充为 $\frg$ 中的一个 $\sl_{2}$ 三元组 $(e,h,f)$。$h$ 在 $\frg$ 上的伴随作用有整数值的权，它把 $\frg$ 分解为权空间 $\frg(n)$，$n\in\ZZ$。令 $\frb=\oplus_{n\geq 0}\frg(n)$。这是 $\frg$ 的一个 Borel 子代数，对应的 Borel 子群就是 $\sB_{e}$ 中唯一的点。

当 $e$ 正则时，Grothendieck 修改的纤维 $\tB_{e}$ 是非约化概形，其底约化概形是单点。$\tB_{e}$ 的坐标环同构于余不变代数 $\Sym(\frt^{*})/(\Sym(\frt^{*})^{W}_{+})$；有趣的是，它也同构于复旗簇 $\calB_{\CC}$ 的 $k$ 系数上同调环。

\subsubsection{} 当 $G=\SL(V)$，$\dim V=n$ 时，$\sB$ 是完备旗 $0=V_{0}\subset V_{1}\subset V_{2}\subset\cdots\subset V_{n-1}\subset V_{n}=V$ 的模空间。Springer 纤维 $\sB_{e}$ 由满足 $eV_{i}\subset V_{i-1}$ 的旗组成。

\subsubsection{}\label{sss:SL3 subreg} 考虑 $G=\SL_{3}$，并在 $V$ 的标准基 $\{v_{1}, v_{2}, v_{3}\}$ 下取 $e=\begin{pmatrix} 0 & 0 & 1\\  0 & 0& 0 \\  0 & 0& 0 \end{pmatrix}$。此时 $\sB_{e}$ 是两个 $\PP^{1}$ 之并：第一个 $\PP^{1}$ 由旗 $0\subset V_{1}\subset \jiao{v_{1}, v_{2}}\subset V$ 组成，其中 $V_{1}$ 在固定平面 $\ker(e)=\jiao{v_{1},v_{2}}$ 内变动；第二个 $\PP^{1}$ 由旗 $0\subset \jiao{v_{1}}\subset V_{2}\subset V$ 组成，其中 $V_{2}$ 是包含 $\Im(e)=\jiao{v_{1}}$ 的变动超平面。

\subsubsection{} 考虑 $G=\SL_{4}$ 及幂零元 $e: v_{3}\mapsto v_{1}\mapsto 0, v_{4}\mapsto v_{2}\mapsto 0$，基取为 $V$ 的标准基 $\{v_{1},v_{2},v_{3},v_{4}\}$。若旗 $0\subset V_{1}\subset V_{2}\subset V_{3}\subset V$ 属于 $\sB_{e}$，则 $V_{1}\subset \ker(e)=\jiao{v_{1},v_{2}}$，且 $V_{3}\supset\Im(e)=\jiao{v_{1},v_{2}}$。记 $H:=\jiao{v_{1},v_{2}}$。分两种情形：
\begin{enumerate}
\item $V_{2}=H$。此时 $V_{1}$ 可取 $H$ 中任一直线，$V_{3}$ 可取任一包含 $V_{2}$ 的超平面。我们得到 $\sB_{e}$ 的一个闭子簇，同构于 $\PP(H)\times \PP(H)\cong \PP^{1}\times \PP^{1}$。把这个闭子簇记为 $C_{1}$。
\item $V_{2}\neq H$。这种情形定义了 $\sB_{e}$ 的一个开子概形 $U$。设直线 $V_{1}\subset H$ 由 $av_{1}+bv_{2}$ 张成，其中 $[a:b]\in\PP(H)$，则 $V_{2}$ 在 $V/H=\jiao{v_{3},v_{4}}$ 中的像由 $av_{3}+bv_{4}$ 张成。固定 $V_{1}$ 后，$V_{2}$ 的选取由 $\Hom(\jiao{av_{3}+bv_{4}}, H/V_{1})\cong\Hom(V_{1},H/V_{1})$ 给出。一旦 $V_{2}$ 固定，$V_{3}=V_{2}+H$ 也随之固定。因此 $U$ 同构于 $\PP(H)\cong \PP^{1}$ 上直线丛 $\calO(2)$ 的全空间。
\end{enumerate}
由上述讨论可见 $\sB_{e}$ 的维数为 $2$，$C_{1}$ 是 $\sB_{e}$ 的不可约分支，$U$ 的闭包也是，记后者为 $C_{2}$。于是 $C_{1}\cup C_{2}=\sB_{e}$，而 $C_{1}\cap C_{2}$ 是 $C_{1}\cong\PP^{1}\times\PP^{1}$ 中的对角线。

\subsubsection{$A$ 型 Springer 纤维的分支}\label{sss:Spr type A} 当 $G=\SL_{n}=\SL(V)$ 时，Spaltenstein \cite{Spa76} 与 Steinberg \cite[]{St} 用 $n$ 阶标准 Young 表给出了 $\sB_{e}$ 的不可约分支的刻画。这与后文讨论的 Springer 对应有关，见 \S\ref{sss:A Spr corr}。下面我们遵循 \cite[Ch II, \S5]{Spaltenstein} 的叙述。

固定幂零元 $e\in\sN$，其 Jordan 型是 $n$ 的分拆 $\l$。这就是说，若分拆写成 $n=\l_{1}+\l_{2}+\cdots$，则 $e$ 的 Jordan 块大小依次为 $\l_{1},\l_{2},\cdots$。我们将构造一个（非代数的）映射 $\sB_{e}\to ST(\l)$，其中 $ST(\l)$ 是分拆 $\l$ 的标准 Young 表构成的离散集。对每个满足 $eV_{i}\subset V_{i-1}$ 的完备旗 $0=V_{0}\subset V_{1}\subset\cdots\subset V_{n-1}\subset V_{n}=V$，$e$ 诱导出 $V/V_{n-i}$ 的幂零自同态。设 $\mu_{i}$ 是 $e$ 在 $V/V_{n-i}$ 上的 Jordan 型，则 $\mu_{i}$ 是 $i$ 的分拆。分拆的递增序列 $\mu_{1},\mu_{2},\cdots, \mu_{n}=\l$ 满足：$\mu_{i}$ 由 $\mu_{i-1}$ 把某一部分加 $1$ 得到（包括新增一个大小为 $1$ 的部分）。由此得到 $\l$ 的 Young 图 $Y(\l)$ 的子图递增序列 $Y_{1},\cdots, Y_{n}=Y(\l)$。把 $Y_{i}-Y_{i-1}$ 中唯一的格子标上 $i$，就得到一张标准 Young 表。

Spaltenstein \cite[Ch II, Prop 5.5]{Spaltenstein} 证明：每张标准 Young 表在 $\sB_{e}$ 中的原像的闭包都是不可约分支。而且 $\sB_{e}$ 的全部不可约分支都这样得到，维数都等于 $d_{e}=\frac{1}{2}\sum_{i}\l^{*}_{i}(\l^{*}_{i}-1)$，其中 $\l^{*}$ 是 $\l$ 的共轭分拆。特别地，最高维上同调 $\cohog{2d_{e}}{\sB_{e}}$ 的维数等于 $\#ST(\l)$，后者也是对称群 $S_{n}$ 的一个不可约表示的维数。这一陈述是 Springer 对应的数值影子：后者断言 $\cohog{2d_{e}}{\sB_{e}}$ 自然地是对称群 $S_{n}$ 的不可约表示。

Spaltenstein \cite[Ch II, Prop 5.9]{Spaltenstein} 还证明：存在 $\sB$ 到仿射空间的分层，使 $\sB_{e}$ 是若干层之并。这说明上同调上的限制映射 $\cohog{*}{\sB}\to \cohog{*}{\sB_{e}}$ 是满射。

\subsubsection{}\label{sss:Sp4} 考虑 $G=\Sp(V)$，其中 $V$ 是 $2n$ 维辛向量空间。此时 $\sB$ 是满足 $V^{\bot}_{i}=V_{2n-i}$（$i=1,\cdots, n$）的完备旗
\begin{equation*}
0=V_{0}\subset V_{1}\subset V_{2}\subset\cdots\subset V_{n}\subset\cdots\subset V_{2n-1}\subset V_{2n}=V
\end{equation*}
的模空间。Springer 纤维 $\sB_{e}$ 由那些对所有 $i$ 满足 $eV_{i}\subset V_{i-1}$ 的旗组成。

考虑 $\dim V=4$ 的情形。选取 $V$ 的基 $\{v_{1},v_{2},v_{3},v_{4}\}$，使 $V$ 上的辛形式 $\om$ 满足：当 $i=1,2$ 时 $\om(v_{i},v_{5-i})=1$，而当 $i+j\neq 5$ 时 $\om(v_{i},v_{j})=0$。设 $e$ 是 $\frg=\sp(V)$ 中由 $e:v_{4}\mapsto v_{1}\mapsto 0, v_{3}\mapsto 0,v_{2}\mapsto 0$ 给出的幂零元。此时 $\sB_{e}$ 中的旗 $0\subset V_{1}\subset V_{2}\subset V_{1}^{\bot}\subset V$ 必须满足 $\jiao{v_{1}}\subset V_{2}\subset\jiao{v_{1},v_{2},v_{3}}$，而这也是它属于 $\sB_{e}$ 的唯一条件（习题 \ref{exI:211Sp4}）。这样的 $V_{2}$ 对应于直线 $V_{2}/\jiao{v_{1}}\subset \jiao{v_{2},v_{3}}$，即 $\PP^{1}=\PP(\jiao{v_{2},v_{3}})$ 中一点。在此 $\PP^{1}$ 之上有万有秩二丛 $\calV_{2}$，它在 $V_{2}/\jiao{v_{1}}$ 处的纤维是二维空间 $V_{2}$。再选取 $V_{1}$，相当于在 $\calV_{2}$ 的射影化中取一点。正合列 $0\to\jiao{v_{1}}\to V_{2}\to V_{2}/\jiao{v_{1}}\to0$ 给出 $\PP^{1}$ 上向量丛的正合列 $0\to\calO\to\calV_{2}\to\calO(-1)\to0$。因此 $\calV_{2}$ 同构于 $\calO(-1)\oplus\calO$，而 $\sB_{e}\cong \PP(\calO(-1)\oplus\calO)$ 是一张 Hirzebruch 曲面。

\subsubsection{次正则 Springer 纤维}\label{sss:subreg} \S\ref{sss:SL3 subreg} 中的例子是{\em 次正则 Springer 纤维}的最简单情形。$\sN$ 中恰有一个余维 $2$ 的幂零轨道 $\calO_{\subreg}$，称为{\em 次正则幂零轨道}。对 $e\in\calO_{\subreg}$，已知 $\sB_{e}$ 是若干 $\PP^{1}$ 之并，其构型如下。可以对 $\sB_{e}$ 作对偶图：顶点为 $\sB_{e}$ 的不可约分支，若两个分支相交（实际上交于单点）则连一条边。

为简单起见设 $G$ 为伴随型。设 $G'$ 是另一个伴随单群，其型如下定义。若 $G$ 单链，取 $G'=G$。若 $G$ 是 $B_{n}$、$C_{n}$、$F_{4}$、$G_{2}$ 型，则分别取 $G'$ 为 $A_{2n-1}$、$D_{n+1}$、$E_{6}$、$D_{4}$ 型。可以证明 $\sB_{e}$ 总是若干 $\PP^{1}$ 之并，其对偶图是 $G'$ 的 Dynkin 图。非单链情形的规则是：每个长单根对应 $2$ 个或 $3$ 个 $\PP^{1}$，每个短单根对应唯一一个 $\PP^{1}$。这样一组 $\PP^{1}$ 称为 Dynkin 曲线，参见 \cite[\S3.10, Definition and Prop 2]{StConj} 与 \cite[\S6.3]{Slodowy}。

例如，当 $G$ 为 $A_{n}$ 型时，$\sB_{e}$ 是由 $n$ 个 $\PP^{1}$ 组成的链：$\sB_{e}=C_{1}\cup C_{2}\cup\cdots\cup C_{n}$，其中 $C_{i}\cap C_{i+1}$ 为一点，其余两两不交。

Brieskorn \cite{Brieskorn} 遵循 Grothendieck 的建议，把幂零锥沿次正则轨道的奇点与 Klein 奇点联系起来，并且在 $\frg$ 内部实现了该奇点的半通用形变。设 $G$ 单链。可以构造过 $e$ 的二维横截片 $S_{e}\subset\sN$，使得 $S_{e}$ 除 $e$ 外都由正则元组成。此时 $S_{e}$ 是正规曲面，在 $e$ 处有 Klein 奇点，其型与 $G$ 的 Dynkin 图相同。\footnote{Klein 奇点指与 $\AA^{2}$ 被 $\SL_{2}$ 的有限子群作商后在 $(0,0)$ 处的奇点解析同构的曲面奇点。} 其原像 $\tilS_{e}:=\pi^{-1}(S_{e})\subset\tN$ 恰是 $S_{e}$ 的极小消解，因此 $\sB_{e}$ 是例外除子之并。把 $\sB_{e}$ 的分支与 $G$ 的单根对应起来后，例外除子的相交矩阵恰为 $G$ 的 Cartan 矩阵的负矩阵。Slodowy \cite{Slodowy} 把上述图景推广到非单链群；关于次正则轨道与 Klein 奇点之间联系的精美叙述，参见他的书 \cite{Slodowy}。
